\Atlas{41a}{4.1 Agentic AI：Agent 应用与基础设施}{4.1 Agentic AI 全景洞察}
\Atlas{41b}{4.1 Agentic AI：模型基础设施}{4.1 Agentic AI 全景洞察}

\BodyPage{4.1 Agentic AI 全景洞察}{任务入口与模型执行同时居前}{ch4-41-position}
两张全景共保留145个项目、25个类别。Agent Infra连接任务入口、构建编排与运行支撑；Model Infra提供模型服务、训练、数据与算力基础。由Codex、OpenClaw等入口继续向下，可以找到Dify、LangChain、vLLM与PyTorch所处的不同功能位置。

2026年6月，个人助手、编码代理和推理引擎的OpenRank分别为1,373.73、1,328.86和1,127.31，合占当月已观测样本的42.3\%。前两类面向任务组织，后一类提供模型执行能力，\key{重点区域同时分布在两张全景中。}
\plate{figures/41/03_category_rank.pdf}{图4-3　全部25类的期末影响力与观测覆盖}{来源：\AgentRef；\href{https://open-digger.cn/docs/user_docs/metrics/global_openrank}{OpenDigger全域OpenRank方法}。2026-06，141/145项有值，合计9,050.38。紫色为Agent Infra，蓝色为Model Infra；条长共用尺度，右列为有值数/名单数。类别值为已观测仓库值之和。}
类别总量也受收录数量影响：编码代理有16项，推理引擎有8项，个人助手有7项。这里的排序用于定位继续调查的区域，不能直接解释为单位项目效率、技术需求或支持优先级。OpenRank反映协作网络影响，不是用户数、模型性能或市场份额。

\BodyPage{4.1 Agentic AI 全景洞察}{月度位置提示变化，也受到观测覆盖影响}{ch4-41-monthly}
个人助手在Agent Infra中从较后位置进入前列；Model Infra的推理引擎则始终位居本层首位。沿全景分类排列的月度矩阵，让任务入口的变化与基础工具的稳定位置可以同时被看见。
\plate{figures/41/04_monthly_positions.pdf}{图4-4　两层全景的月度类别位置}{来源：\AgentRef。2025-07至2026-06。数字为各自层内的类别名次，颜色越深越靠前；并列按同名次处理，无观测留空。Agent与Model分别计算，名次不直接跨层比较。}
\key{位次变化不等于同一批项目的增长。}个人助手7项中，只有2项在整个主期连续有值；类别位置的移动也包含后来进入观测的项目。编码工作流与工具链在2025年7月没有观测，空白不能解释为相关能力尚不存在。

因此，既要保留当前全景来发现新的可见对象，也要用固定的连续样本观察既有项目。两种范围回答不同的问题，需要在同一叙事中并读。

\BodyPage{4.1 Agentic AI 全景洞察}{观测范围改变，趋势方向也可能改变}{ch4-41-trajectories}
期末前三类的曲线并不相同：个人助手在主期后半段抬升较明显，编码代理与推理引擎呈现另一种走势。分面共用刻度后，差异更容易辨认；进一步固定对象，整体轨迹又出现变化。
\plate{figures/41/05_category_trends.pdf}{图4-5　重点功能区的分值轨迹与观测范围}{来源：\AgentRef。2025-07至2026-06。上图为当月有值项目合计，三类共用0—4,000刻度；下图比较全名单已观测合计与107项全期共同样本，两线共用纵轴。上行数字为各月有值项目数；缺失不补零。}
2026年1月至3月，全名单已观测合计由13,916.42升至15,009.23，107项共同样本却由12,705.10降至11,486.11。期末26.6\%的合计来自非全期连续有值项目，不能把全部变化归给同一批既有项目。

\key{覆盖扩展与既有项目变化应分别观察。}这里描述冻结指标的轨迹；历史算法版本与跨月归一化参数未完全固定，不能据共同回落推断产业衰退，也不能据覆盖增加推断全球参与增长。

\BodyPage{4.1 Agentic AI 全景洞察}{份额上升，不一定意味着绝对扩张}{ch4-41-share}
在连续有值的107项中，服务部署、编码代理、多代理编排以及编译器与加速的份额上升较多。这些类别分布于不同层级，相对变化并不集中于单一技术路线。
\plate{figures/41/08_common_shift.pdf}{图4-6　共同样本在全景各类别中的份额变化}{来源：\AgentRef。固定107项，比较2025-07与2026-06端点；分母为共同样本当月OpenRank合计。紫色为Agent，蓝色为Model，右列为共同项数；无共同样本不解释为稳定。}
服务部署的份额由2.83\%升至6.27\%，分值由412.26变为416.59；编码代理份额由14.83\%升至17.59\%，分值却由2,161.89降至1,169.58。\key{同样是份额上升，一类分值基本稳定，另一类分值下降。}

共同样本减少了覆盖差异，也排除了中途进入和缺少记录的项目。例如个人助手仅2项进入共同样本，不能据此替代其当前7项全貌。相对位置、绝对轨迹与当前名单需要同时保留。

\BodyPage{4.1 Agentic AI 全景洞察}{从功能区进入项目，辨认不同的集中结构}{ch4-41-projects}
推理引擎8项在期末均有值；vLLM、SGLang与TensorRT-LLM在主期每月都居前三，期末合占该类的66.7\%。头部组合稳定，分值与相互差距仍在变化。
\plate{figures/41/06_inference_projects.pdf}{图4-7　推理引擎的项目位置与完整月度记录}{来源：\AgentRef。按2026-06值排列；轨迹为2025-07至2026-06，每行统一0—800刻度。缺失月份断线，不补零；相关仓库不视为完全独立生态。}
个人助手更集中：OpenClaw与Hermes Agent期末分别为592.19和518.86，两项合占该类的80.9\%。同样进入期末前三类的功能区，内部结构并不相同。
\plate{figures/41/07_within_categories.pdf}{图4-8　个人助手完整构成与三类头部集中度}{来源：\AgentRef。2026-06，均以各类别有值项目合计为分母。前三项占比分别为85.3\%、67.3\%、66.7\%，不是用户或技术依赖份额。}
\key{完整全景定位能力，类内对照定位对象。}集中度有助于选择后续维护调查的入口，不能单凭指标判断系统脆弱性、市场支配或资助顺序。

\Atlas{42}{4.2 AI 大模型：任务全景}{4.2 AI 大模型全景洞察}

\BodyPage{4.2 AI 大模型全景洞察}{模型条目分散，不等于发布主体分散}{ch4-42-structure}
全景按六类任务各取近期下载Top 16，保留96个模型仓库。嵌入与检索的第一项占本类41.2\%，前三项合占66.6\%；多模态区相应为14.3\%和35.2\%。同样16项，内部集中结构并不相同。
\plate{figures/42/02_concentration.pdf}{图4-10　六类近期下载的样本内集中程度}{来源：\DataRef，2026-10-02期后快照。分母为各类近期下载Top 16的下载合计；空心点为第一项，实心点为前三项。不是全领域集中度或市场份额。}
按发布账户归组，多模态16项来自4个账户；Qwen与google合占80.3\%的样本内下载。\key{仓库层面较分散的结构，可以同时集中于少数发布账户。}
\plate{figures/42/04_publishers.pdf}{图4-11　两类样本的发布账户构成}{来源：\DataRef，同一期后快照。按完全一致的publisher字段汇总；右列为下载份额与仓库条数。账户未经机构合并，不等于最终母公司或全部贡献者。}
\src{\href{https://huggingface.co/docs/hub/models-download-stats}{Hugging Face下载计数}基于指定文件请求，不能解释为独立用户或部署数；不同模型库与格式的计数规则也可能不同。}

\BodyPage{4.2 AI 大模型全景洞察}{保留完整类内对照，区分模型、版本与账户}{ch4-42-detail}
检索区同时出现向量表示与相关性排序模型；多模态区包含Qwen、Gemma等不同版本。以下保留两类全部32条，避免由头部几个名称代替完整结构。
\plate{figures/42/03_full_category_detail.pdf}{图4-12　检索与多模态的完整类内下载结构}{来源：\DataRef，2026-10-02期后快照。每列各16项，以各自下载合计为100\%；条长共用0—45\%尺度，行序按原表名次。完整模型标识见全景及随附数据。}
多模态区lmstudio-community的4条记录属于同一命名家族的MLX 4bit、5bit、6bit与8bit版本。将它们视为四种独立能力会夸大多样性，直接合并又会丢失运行格式的维护对象。\key{本节保留原仓库，家族与派生关系另行核对。}

\BodyPage{4.2 AI 大模型全景洞察}{改变观察指标，改变的往往是入选对象}{ch4-42-overlap}
近期下载反映近30天的文件获取信号，累计下载保留更长历史的积累，点赞记录关注。三项指标按六类任务分别取Top 16，共288条榜单记录，对应194个唯一模型仓库，三榜同时出现的仅22个。
\plate{figures/42/05_overlap.pdf}{图4-13　改变观察指标后保留了多少共同模型}{来源：\DataRef，同一原表快照（2026-10-02，期后）。每个彩色小方块为1个共同模型，灰色为未重合名额；分母为每榜每类固定的16个名额。各列为两榜交集，不是增长率。}
近期与累计下载共有64个模型，近期下载与点赞仅共有27个。图像与视频区的近期下载与点赞榜只有2个共同条目，累计下载与点赞榜只有1个。

该区近期下载首位为timm的MobileNetV3 Small，累计下载首位为Falconsai的nsfw\_image\_detection，点赞首位则为black-forest-labs的FLUX.1-dev。\key{频繁获取与受到关注，把读者带向不同任务和对象。}

三种视图来自同一快照，不能顺次连接成时间曲线；没有进入某一Top 16也不能补成零值。持续观察需要先固定要回答的问题，再固定对象、任务边界和时间窗口。

\BodyPage{4.2 AI 大模型全景洞察}{机器人样本：近期获取、历史积累与关注各有重点}{ch4-42-robotics}
机器人区把三种观察的差异呈现得更清楚。近期下载首位GR00T-N1.7-3B占本类18.9\%；累计下载首位OpenVLA占84.1\%；点赞首位SmolVLA占12.2\%。这些比例分别来自三份独立选样的榜单。
\plate{figures/42/06_robotics.pdf}{图4-14　机器人模型在三种观察中的头部构成}{来源：\DataRef，2026-10-02期后快照。每榜展示前5项，分母仍为对应榜单全部16项；不是将前5项重新归一化。三榜分别选样，下载与点赞不等于设备部署、使用者或执行效果。}
\key{同一任务区可以同时保留探索热点、历史积累与社区关注。}机器人模型能否用于实际任务，还要确认观测输入、动作表示、控制器与物理设备之间的匹配。下载信号能够提示调查对象，不能替代运行环境中的验证。

\BodyPage{4.2 AI 大模型全景洞察}{放宽视野：专用模型与运行工具同样重要}{ch4-42-hf}
HF夏季报告显示，在2026年7月29日的累计下载分布中，参数少于10亿的模型占全Hub下载的61.5\%；约26\%对应缺少参数规模信息的仓库。只在已标注参数的部分归一化，小模型比例才约为83\%。
\plate{figures/42/08_hf_parameter.pdf}{图4-15　模型参数规模与累计下载分布}{来源：\HFRef，2026-07-29期后快照。按官方图精确标签复绘，未取得逐模型底表；分母含无参数元数据仓库，原图舍入合计99.9\%，未补造尾差。}
\key{技术观察也要覆盖可嵌入具体工作流的专用能力。}这一结果与全景中的MiniLM、MobileNet等对象相呼应，但不能据下载份额判断小模型在所有任务中更有效或更节能。
\plate{figures/42/09_hf_runtime.pdf}{图4-16　模型仓库声明的运行库与格式增长}{来源：\HFRef，2026-01-07至2026-07-29，跨越主期截止日。按官方印刷标签复绘；统计声明相应library的模型仓库数，不是下载、安装或用户增长。标签未必互斥，不能相加为Hub总量。}
GGUF的高增幅来自较小基数，Transformers的绝对仓库数仍大得多。MLX、LeRobot等变化则让运行适配和具身工具进入视野：模型能否被采用，还取决于格式、工具与接口。这是更宽的生态背景，不是对前述96条样本变化的因果解释。

\BodyPage{4.2 AI 大模型全景洞察}{回到任务，检查能力之间能否接续}{ch4-42-task}
全景帮助寻找模型，模型卡进一步说明其输入、输出和使用条件。若要把多个模型用于同一任务，还需要逐环节核对。下图归纳可按任务组合的能力，并不表示这些项目已被组合部署。
\plate{figures/42/07_task_paths.pdf}{图4-17　从任务链定位模型与检验环节}{来源：全景分类及相关模型卡的阅读归纳，非部署实测。资料问答与机器人行示意任务环节；语音行列出可组合能力，无需全部采用或固定串行。}
\topic{资料问答：分别检验表示、检索与回答}
\href{https://huggingface.co/sentence-transformers/all-MiniLM-L6-v2}{all-MiniLM-L6-v2}将句段转换为向量，cross-encoder为查询与候选资料计算相关性；检索能否支撑回答，还需要引用核对与生成环节的检验。
\topic{语音服务：区分识别、说话人分析与生成}
\href{https://huggingface.co/openai/whisper-small}{Whisper-small}识别语音，\href{https://huggingface.co/hexgrad/Kokoro-82M/blob/f3ff3571791e39611d31c381e3a41a3af07b4987/README.md}{Kokoro-82M}合成语音，\href{https://huggingface.co/pyannote/speaker-diarization-3.1/blob/84fd25912480287da0247647c3d2b4853cb3ee5d/README.md}{pyannote}区分谁在何时说话。三者输出和评测任务不同，采用前还要检查语言、噪声与说话人条件。
\topic{机器人任务：验证动作接口与设备条件}
\href{https://huggingface.co/openvla/openvla-7b}{OpenVLA}提供视觉语言动作模型与推理说明；实际部署仍须确认动作归一化、机器人接口与安全约束。功能位置、下载名次都不能替代具体场景中的验证。

\Atlas{43}{4.3 云基础设施全景}{4.3 云基础设施全景}
\Atlas{44}{4.4 数据库系统全景}{4.4 数据库系统全景}
\Atlas{45}{4.5 RISC-V 技术全景}{4.5 RISC-V 技术全景}
\Atlas{46}{4.6 PyTorch 生态全景}{4.6 PyTorch 生态全景}
